\documentclass[acmsmall,review,screen]{acmart}

%% --- ACM metadata --------------------------------------------------------
\acmJournal{PACMPL}
\acmVolume{1}
\acmNumber{HASKELL}
\acmYear{2026}
\acmDOI{}

%% --- Packages ------------------------------------------------------------
\usepackage{listings}
\usepackage{amsmath}
\let\Bbbk\relax
\usepackage{amssymb}
\usepackage{stmaryrd}   % \llbracket, \rrbracket
\usepackage{mathtools}
\usepackage{booktabs}
\usepackage{microtype}

%% --- Haskell listings style ----------------------------------------------
\lstdefinelanguage{Haskell}{
  keywords={data, type, class, instance, where, let, in, do, if, then, else,
            case, of, import, module, deriving, infixl, infixr},
  keywordstyle=\color{blue!70!black}\bfseries,
  sensitive=true,
  comment=[l]{--},
  morecomment=[s]{\{-}{-\}},
  commentstyle=\color{gray}\itshape,
  stringstyle=\color{orange!80!black},
  morestring=[b]",
  literate=
    {<>}{{$\diamond$}}1
    {/\\}{{$\wedge$}}1
    {\\\\}{{$\setminus$}}1
    {->}{{$\to$}}2
    {=>}{{$\Rightarrow$}}2
    {forall}{{$\forall$}}1
    {>>}{{$>\!\!>$}}2
    {>>=}{{$>\!\!>\!\!=$}}3,
}
\lstset{
  language=Haskell,
  basicstyle=\ttfamily\small,
  columns=flexible,
  keepspaces=true,
  frame=none,
  xleftmargin=1em,
}

%% --- Math macros ---------------------------------------------------------
\newcommand{\RE}{\mathsf{RE}}
\newcommand{\Eff}{\mathsf{Effectful}}
\newcommand{\pre}{\mathit{pre}}
\newcommand{\post}{\mathit{post}}
\newcommand{\fut}{\mathit{future}}
\newcommand{\ret}{\mathit{ret}}
\newcommand{\anything}{\mathtt{anything}}
\newcommand{\finally}[1]{\mathtt{finally}(#1)}
\newcommand{\derives}[2]{\partial_{#1}(#2)}
\newcommand{\norm}[1]{\lceil #1 \rceil}
\newcommand{\sem}[1]{\llbracket #1 \rrbracket}

\begin{document}

%% --- Title / Authors -----------------------------------------------------
\title{Effectful Computations with Future Conditions:\\
       A Monadic Approach to Temporal Specification}

\subtitle{Functional Pearl}

\author{Author Name}
\affiliation{%
  \institution{Institution}
  \country{Country}
}
\email{author@institution.edu}

\begin{abstract}
We present \emph{FutureCond}, a Haskell library that extends the classical
pre/post-condition contract model with \emph{future conditions}: regular
expressions over events that specify what the remainder of a program must
eventually do. Future conditions are carried as a first-class field in the
\lstinline{Effectful} monad. When operations are composed via monadic bind,
future conditions are propagated and updated automatically using a
Brzozowski-derivative-based subtraction operator. A future condition that
normalises to $(\texttt{\_})^*$ certifies that all temporal obligations
introduced by the program have been discharged. We demonstrate the approach
across seven domains---memory management, file handles, mutexes, database
transactions, cryptographic sessions, network protocols, and IoT sensor
control---showing that resource-lifecycle and ordering properties that are
invisible to pre/post reasoning become statically visible as residual future
conditions.
\end{abstract}

\keywords{temporal specification, monadic effects, regular expressions,
          Brzozowski derivatives, resource lifecycle, Haskell}

%%
%% ACM Classification
%%
\begin{CCSXML}
<ccs2012>
  <concept>
    <concept_id>10011007.10011006.10011008.10011009.10011012</concept_id>
    <concept_desc>Software and its engineering~Functional languages</concept_desc>
    <concept_significance>500</concept_significance>
  </concept>
  <concept>
    <concept_id>10011007.10011006.10011008.10011024.10011028</concept_id>
    <concept_desc>Software and its engineering~Data types and structures</concept_desc>
    <concept_significance>300</concept_significance>
  </concept>
  <concept>
    <concept_id>10011007.10011006.10011041</concept_id>
    <concept_desc>Software and its engineering~Software verification and validation</concept_desc>
    <concept_significance>300</concept_significance>
  </concept>
</ccs2012>
\end{CCSXML}
\ccsdesc[500]{Software and its engineering~Functional languages}
\ccsdesc[300]{Software and its engineering~Data types and structures}
\ccsdesc[300]{Software and its engineering~Software verification and validation}

\maketitle

%% =========================================================================
\section{Introduction}
\label{sec:intro}
%% =========================================================================

Pre- and post-conditions are the workhorse of modular program specification.
A pre-condition constrains the state in which an operation may be called;
a post-condition describes what holds once it returns. Together they reason
about a single \emph{call boundary}. Yet many correctness properties of
real programs span multiple calls separated by an unbounded amount of
intervening computation:

\begin{itemize}
  \item Every \lstinline{malloc(a)} must eventually be followed by
        \lstinline{free(a)}.
  \item Every \lstinline{open(f)} must eventually be followed by
        \lstinline{close(f)}.
  \item Every \lstinline{beginTx} must eventually reach a
        \lstinline{commit} or \lstinline{rollback}.
  \item Every \lstinline{acquire(m)} must eventually be followed by
        \lstinline{release(m)}.
\end{itemize}

These are \emph{deferred obligations}---temporal constraints on what the
\emph{rest} of the program must do after an operation is invoked.
Pre/post reasoning cannot express them: a post-condition records what is
true \emph{now}, not what must become true \emph{later}.

We propose \emph{future conditions} as a lightweight, compositional
mechanism for specifying such obligations. A future condition is a regular
expression over events that the remaining execution trace must match. It is
carried as a third annotation alongside \lstinline{pre} and \lstinline{post}
in an \lstinline{Effectful} monad. Crucially, when two effectful operations
are sequenced, future conditions are composed automatically: the bind
operator uses a Brzozowski-style derivative to remove from the first
operation's future condition those obligations that the second operation
will satisfy, and then intersects the residual with the second operation's
own future condition.

The key insight is that the universal future condition---``anything may
happen''---acts as the \emph{identity}: an operation whose future normalises
to $(\mathtt{\_})^*$ has no outstanding obligations. The normalised future
of a composed program therefore serves as a certificate of temporal
correctness, or as a precise diagnosis of which obligations remain open.

\paragraph{Contributions.}
\begin{enumerate}
  \item We define the \lstinline{Effectful} monad (\S\ref{sec:effectful}),
        parameterised over a \lstinline{Composable} algebra, with
        \lstinline{pre}, \lstinline{post}, and \lstinline{future} fields,
        and prove that it satisfies the monad laws
        (\S\ref{sec:monad-laws}).
  \item We instantiate the algebra to regular expressions over events,
        using Brzozowski derivatives to define the subtraction operator that
        drives future-condition propagation (\S\ref{sec:re}).
  \item We demonstrate the approach on seven real-world use cases
        (\S\ref{sec:examples}), showing that resource leaks and protocol
        violations manifest as non-trivial residual future conditions.
\end{enumerate}

%% =========================================================================
\section{Background}
\label{sec:background}
%% =========================================================================

\subsection{Regular Expressions over Events}

Events are pairs $(\mathit{name}, \mathit{args})$ where $\mathit{args}$
is a list of \emph{terms}---variables, string literals, or integer literals.
The regular expression type $\RE$ is:

\[
  r \;::=\; \bot \mid \varepsilon \mid e \mid \_ \mid r_1 \cdot r_2
          \mid r_1 \vee r_2 \mid r_1 \wedge r_2 \mid r^*
\]

where $\bot$ is the empty language, $\varepsilon$ the empty trace,
$e$ matches a single specific event, $\_$ (wildcard) matches any single
event, $\cdot$ is concatenation, $\vee$ union, $\wedge$ intersection,
and $^*$ Kleene star.

In Haskell:

\begin{lstlisting}
data RE = Bot | Epsilon | Single Event | Wildcard
        | Seq RE RE | Or RE RE | And RE RE | Star RE
\end{lstlisting}

Two derived forms are central to temporal specification:

\begin{align*}
  \anything &\triangleq (\mathtt{\_})^* \\
  \finally{e} &\triangleq \anything \cdot e \cdot \anything
\end{align*}

$\anything$ matches any trace (no obligation); $\finally{e}$ requires
event $e$ to appear somewhere in the remaining trace.

\subsection{Brzozowski Derivatives}

The \emph{Brzozowski derivative}~\cite{brzozowski1964} of a language $L$
with respect to a word $w$ is
$\partial_w(L) = \{ v \mid wv \in L \}$.
For regular expressions, derivatives can be computed symbolically.
We use a first-set-based variant: \lstinline{first r} returns the set of
\emph{first symbols} of $r$ as \lstinline{Fst} values (\lstinline{Pos e},
\lstinline{Neg e}, or \lstinline{W}), and \lstinline{derivitive e r}
computes $\partial_e(r)$.

The subtraction operator $r_1 \setminus r_2$---``what of $r_2$ remains to
be satisfied after observing the traces described by $r_1$''---is defined
by unfolding $r_1$ event by event using its first set:

\[
  \varepsilon \setminus r_2 = r_2
  \qquad
  r_1 \setminus r_2 = \bigvee_{e \,\in\, \mathit{first}(r_1)}
                        \bigl(\norm{r_1'} \setminus \norm{r_2'}\bigr)
\]

where $r_1' = \derives{e}{r_1}$, $r_2' = \derives{e}{r_2}$, and
$\norm{\cdot}$ denotes normalisation.

%% =========================================================================
\section{The \texttt{Effectful} Monad}
\label{sec:effectful}
%% =========================================================================

\subsection{The \texttt{Composable} Class}

We abstract over the specification algebra via a type class:

\begin{lstlisting}
class Composable a where
  concatenation :: a -> a -> a
  conjunction   :: a -> a -> a
  empty         :: a
  universe      :: a
  subtraction   :: a -> a -> a
  normalize     :: a -> a
\end{lstlisting}

Operators \lstinline{<>} (concatenation), \lstinline{/\\} (conjunction),
and \lstinline{\\\\} (subtraction) are defined as infix synonyms.
\lstinline{empty} is the identity for postcondition sequencing ($\varepsilon$)
and \lstinline{universe} is the identity for future conditions ($(\mathtt{\_})^*$).

\subsection{The \texttt{Effectful} Type}

\begin{lstlisting}
data Effectful eff a = Effectful
  { ret    :: a
  , pre    :: eff
  , post   :: eff
  , future :: eff
  }
\end{lstlisting}

An \lstinline{Effectful eff a} value represents a computation that:
\begin{itemize}
  \item returns a value of type \lstinline{a},
  \item requires \lstinline{pre} to have been observed in the trace so far,
  \item produces \lstinline{post} as its observable effect, and
  \item requires the \emph{remaining} trace to satisfy \lstinline{future}.
\end{itemize}

\subsection{Functor, Applicative, and Monad Instances}

\lstinline{fmap} maps over the return value without altering specifications:

\begin{lstlisting}
instance Functor (Effectful eff) where
  fmap f e = e { ret = f (ret e) }
\end{lstlisting}

\lstinline{pure} introduces a computation with no effect and no obligation:

\begin{lstlisting}
pure x = Effectful
  { ret = x, pre = universe, post = empty, future = universe }
\end{lstlisting}

The bind operator sequences two computations and propagates future
conditions:

\begin{lstlisting}
e >>= f = Effectful
  { ret    = ret (f (ret e))
  , pre    = pre e <> pre (f (ret e))
  , post   = post e <> post (f (ret e))
  , future = (post (f (ret e)) \\ future e)
             /\ future (f (ret e))
  }
\end{lstlisting}

The future condition of $e \mathbin{>\!\!>\!\!=} f$ deserves close attention.
Let $e_2 = f(\ret(e))$.

\begin{equation}
  \fut(e \mathbin{>\!\!>\!\!=} f)
  \;=\;
  \bigl(\post(e_2) \setminus \fut(e)\bigr)
  \;\wedge\;
  \fut(e_2)
  \label{eq:future-bind}
\end{equation}

The first conjunct removes from $e$'s future obligation those traces that
$e_2$'s post-condition already provides. The second conjunct adds $e_2$'s
own future obligation. The conjunction means \emph{both} residuals must
hold.

\subsection{Intuition via an Example}

Consider:

\begin{lstlisting}
malloc :: Int -> Effectful RE ()
malloc a = Effectful
  { ret = (), pre = universe
  , post   = Single ("malloc", [Num a])
  , future = finally ("free", [Num a]) }

free :: Int -> Effectful RE ()
free a = Effectful
  { ret = (), pre = universe
  , post   = Single ("free", [Num a])
  , future = anything }
\end{lstlisting}

After \lstinline{malloc 1}, the future condition is
$\finally{\mathtt{free}(1)}$. When we then bind \lstinline{free 1}, whose
post-condition is $\mathtt{free}(1)$:

\[
  \mathtt{free}(1) \setminus \finally{\mathtt{free}(1)}
  \;=\; \anything
\]

The subtraction discharges the obligation. The conjunction with
\lstinline{free}'s own future ($\anything$) leaves $\anything$---no
outstanding obligations.

If instead we bind \lstinline{free 2}, the post-condition $\mathtt{free}(2)$
does not match the obligation $\finally{\mathtt{free}(1)}$, and the
residual future contains a non-trivial \lstinline{finally} term indicating
the leak.

%% =========================================================================
\section{The RE Instance}
\label{sec:re}
%% =========================================================================

The \lstinline{Composable RE} instance maps library operations directly onto
\lstinline{RE} constructors:

\begin{lstlisting}
instance Composable RE where
  concatenation = Seq
  conjunction   = And
  empty         = Epsilon
  universe      = anything
  subtraction Epsilon r2 = r2
  subtraction r1 r2 =
    let fs  = first r1
        res = map (\e -> subtraction
                           (normalize (derivitive e r1))
                           (normalize (derivitive e r2))) fs
    in foldr Or Bot res
  normalize = {- algebraic simplification -} ...
\end{lstlisting}

The \lstinline{normalize} function applies algebraic laws to eliminate
redundant constructors: $\varepsilon \cdot r = r$, $\bot \vee r = r$,
$r \wedge (\mathtt{\_})^* = r$, and so on. Normalisation is essential both
for readability of results and for the equality check that terminates the
subtraction recursion.

%% =========================================================================
\section{Examples}
\label{sec:examples}
%% =========================================================================

We demonstrate FutureCond across seven domains. In each case we define
primitive operations, compose them into programs, and inspect the normalised
future condition of the result.

\subsection{Memory Management}

The canonical example from the library's \lstinline{Main.hs}:

\begin{lstlisting}
leak :: Effectful RE ()
leak = do { malloc 1; malloc 2; free 1 }
\end{lstlisting}

\lstinline{normalize (future leak)} evaluates to
\lstinline{finally("free", [2])}, naming the leaked address precisely.

\subsection{File Handle Lifecycle}

\begin{lstlisting}
fh_bad :: Effectful RE ()
fh_bad = do { fileOpen "a.txt"; fileOpen "b.txt"; fileClose "a.txt" }
\end{lstlisting}

Residual: \lstinline{finally("close", ["b.txt"])}.

\subsection{Mutex / Lock Lifecycle}

\begin{lstlisting}
deadlock :: Effectful RE ()
deadlock = do { acquire 1; acquire 2; work; release 1 }
\end{lstlisting}

Residual: \lstinline{finally("release", [2])}.

\subsection{Database Transactions}

\lstinline{beginTx} carries a disjunctive future condition:

\begin{lstlisting}
future = Or (finally ("commit", [])) (finally ("rollback", []))
\end{lstlisting}

A program that calls \lstinline{beginTx} and then \lstinline{commit}
discharges the obligation; one that ends without either leaves the full
disjunction outstanding.

\subsection{Cryptographic Sessions and Nonces}

\begin{lstlisting}
nonceLeak :: Effectful RE ()
nonceLeak = do
  initSession "s"; generateNonce 42
  encrypt "s" "msg"; finalizeSession "s"
\end{lstlisting}

Residual: \lstinline{finally("consumeNonce", [42])}, flagging the
use-once nonce that was never consumed, which could enable a replay attack.

\subsection{Network Protocol (TCP-like)}

\begin{lstlisting}
stalledHandshake :: Effectful RE ()
stalledHandshake = do { sendSYN }
\end{lstlisting}

Residual: \lstinline{finally("recvSYNACK", [])}.

\begin{lstlisting}
teardownMissed :: Effectful RE ()
teardownMissed = do
  sendSYN; recvSYNACK; sendACK; sendData "payload"; sendFIN
\end{lstlisting}

Residual: \lstinline{finally("recvFINACK", [])}.

\subsection{IoT Sensor and Actuator Control}

\begin{lstlisting}
motorLeftRunning :: Effectful RE ()
motorLeftRunning = do { motorOn 3; actuate "fan" 50 }
\end{lstlisting}

Residual: \lstinline{finally("motorOff", [3])}, indicating a safety
hazard---the motor is never shut down.

\paragraph{Summary.}
Table~\ref{tab:examples} summarises the examples.

\begin{table}[h]
\caption{Example programs and their residual future conditions.}
\label{tab:examples}
\begin{tabular}{lll}
\toprule
Domain & Bad program & Residual future \\
\midrule
Memory      & \lstinline{malloc 1,2; free 1}   & \lstinline{finally(free(2))} \\
File handle & \lstinline{open a,b; close a}    & \lstinline{finally(close("b"))} \\
Mutex       & \lstinline{acq 1,2; rel 1}       & \lstinline{finally(release(2))} \\
Transaction & \lstinline{beginTx; write}        & \lstinline{finally(commit) \(\vee\) finally(rollback)} \\
Crypto      & nonce never consumed             & \lstinline{finally(consumeNonce(42))} \\
Network     & FIN without FIN-ACK              & \lstinline{finally(recvFINACK)} \\
IoT         & motor never off                  & \lstinline{finally(motorOff(3))} \\
\bottomrule
\end{tabular}
\end{table}

%% =========================================================================
\section{Monad Laws}
\label{sec:monad-laws}
%% =========================================================================

We verify the monad laws for \lstinline{Effectful} under the assumption that
the \lstinline{Composable} instance is a monoid for concatenation
(\lstinline{empty} is the unit, \lstinline{<>} is associative) and that
\lstinline{universe} is the unit for \lstinline{/\\}.

\paragraph{Left identity} \lstinline{pure x >>= f = f x}.
\begin{align*}
  \pre(\mathtt{pure}\;x \mathbin{>\!\!>\!\!=} f)
    &= \mathit{universe} \mathbin{<>} \pre(f\,x) = \pre(f\,x) \\
  \post(\mathtt{pure}\;x \mathbin{>\!\!>\!\!=} f)
    &= \mathit{empty} \mathbin{<>} \post(f\,x) = \post(f\,x) \\
  \fut(\mathtt{pure}\;x \mathbin{>\!\!>\!\!=} f)
    &= (\post(f\,x) \setminus \mathit{universe})
       \wedge \fut(f\,x)
     = \anything \wedge \fut(f\,x)
     = \fut(f\,x)
\end{align*}

\paragraph{Right identity} \lstinline{e >>= pure = e}.
\begin{align*}
  \pre(e \mathbin{>\!\!>\!\!=} \mathtt{pure})
    &= \pre(e) \mathbin{<>} \mathit{universe} = \pre(e) \\
  \post(e \mathbin{>\!\!>\!\!=} \mathtt{pure})
    &= \post(e) \mathbin{<>} \mathit{empty} = \post(e) \\
  \fut(e \mathbin{>\!\!>\!\!=} \mathtt{pure})
    &= (\mathit{empty} \setminus \fut(e)) \wedge \mathit{universe}
     = \fut(e) \wedge \anything
     = \fut(e)
\end{align*}

\paragraph{Associativity} follows from associativity of \lstinline{<>}
and the distributivity of \lstinline{\\} over \lstinline{/\\}.

%% =========================================================================
\section{Related Work}
\label{sec:related}
%% =========================================================================

\paragraph{Future conditions.}
Song et al.\ propose future conditions as a first-class extension to
pre/post-condition contracts~\cite{song2025futurecond}, formalising the
notion that a function may leave temporal obligations for its caller to
discharge. FutureCond realises this idea as a Haskell library, providing
a monadic interface and derivative-based propagation that the formal
system does not address.

\paragraph{Temporal property verification and repair.}
ProveNFix~\cite{song2024provenfix} uses temporal properties expressed as
regular expressions to guide automated program repair, verifying that
patches satisfy the specification. FutureCond is complementary: rather
than repairing programs post-hoc, it carries temporal obligations as
first-class values so that violations are surfaced statically.

\paragraph{Effect systems and monads.}
The \lstinline{Effectful} monad is related to effect handlers~\cite{pretnar2015}
and graded/indexed monads~\cite{orchard2014,katsumata2014} that track
effects in types. Our contribution is orthogonal: we focus on the
\emph{temporal} dimension of effects---what must happen \emph{after}---rather
than on tracking which effects are present.

\paragraph{Session types.}
Session types~\cite{honda1993,gay2010} enforce communication protocols
at the type level, including ordering constraints. FutureCond is more
lightweight: it operates at the value level, uses regular expressions
rather than types, and can be applied to any domain beyond communication.

\paragraph{Runtime monitoring and LTL.}
Runtime verification~\cite{leucker2009} checks temporal properties of
execution traces against LTL or regular-language specifications.
FutureCond is a \emph{static} approach: the future condition of a program
is computed symbolically before any execution takes place, acting as a
whole-program certificate.

\paragraph{Typestate.}
Typestate systems~\cite{strom1986,bierhoff2007} track resource lifecycle
in types, requiring type-system changes. FutureCond is a library-level
approach requiring no language extensions beyond standard Haskell type
classes.

\paragraph{Brzozowski derivatives.}
Derivatives of regular expressions were introduced by
Brzozowski~\cite{brzozowski1964} and have found wide application in
parsing~\cite{owens2009,might2011} and model checking. Our subtraction
operator is a novel application of the derivative to future-condition
propagation in a monadic setting.

%% =========================================================================
\section{Discussion and Future Work}
\label{sec:discussion}
%% =========================================================================

\paragraph{Decidability.}
Equality of normalised regular expressions is decidable, so the fixed-point
computation in \lstinline{subtraction} terminates whenever the derivative
sequence is finite---which holds for the regular languages used in our
examples. Extending to context-free or $\omega$-regular specifications
remains future work.

\paragraph{Pre-conditions.}
The current implementation tracks \lstinline{pre} as a trace
specification, but does not yet check that \lstinline{pre} is satisfied
by the \lstinline{post} of preceding operations. Integrating pre-condition
verification with the derivative machinery would yield a fully automated
Hoare-logic-style checker within the monad.

\paragraph{Integration with effect handlers.}
An interesting direction is lifting FutureCond into an effect-handler
framework (e.g., \texttt{effectful} or \texttt{polysemy}), so that
temporal specifications can be composed alongside algebraic effect rows.

\paragraph{Tool support.}
The normalised future condition is a human-readable diagnostic, but
tooling to visualise the automaton structure of residual conditions and
relate it back to source locations would improve usability.

%% =========================================================================
\section{Conclusion}
\label{sec:conclusion}
%% =========================================================================

We have presented FutureCond, a small but expressive Haskell library that
extends effectful computations with future conditions---regular expressions
over events that specify temporal obligations on the rest of the program.
The key technical contribution is the bind operator (Equation~\ref{eq:future-bind}),
which uses Brzozowski-derivative-based subtraction to propagate and update
future conditions automatically as operations are sequenced. The resulting
system is lightweight (a single module, no language extensions), composable
(a standard \lstinline{Monad} instance), and informative (residual future
conditions name unmet obligations precisely). We believe future conditions
are a natural and useful complement to pre/post contracts in the Haskell
programmer's specification toolkit.

%% =========================================================================
\bibliographystyle{ACM-Reference-Format}
\bibliography{references}

\end{document}
